%=======================================================================
%  CSE558 Data Science - Midterm Project Report Template
%  Monsoon 2026, IIIT-Delhi
%
%  How to use this template
%   1. Set your project title in the EDIT THIS block below.
%   2. Replace everything in LIGHT BLUE with your own text and numbers.
%      The numbers in the example are made up for illustration.
%   3. Grey italic lines are instructions. Delete them as you write, or
%      hide them all at once with \showhintsfalse.
%   4. Upload your plots to the project and put their file names in
%      \figfile{...}{...}. Until a file exists, a labelled box is shown.
%
%  Sections 1-5 follow items 1-5 of the report guidelines; Section 6 is
%  the limitations of your analysis. The link to your code (item 6) goes
%  at the top. At most 8 pages.
%  Compile with pdfLaTeX (Overleaf's default setting works), twice.
%=======================================================================
\documentclass[10pt,a4paper,twocolumn]{article}

%========================= EDIT THIS =========================
% Your project title: it sets the heading AND the PDF's title.
\newcommand{\projecttitle}{Customer Behaviour and Repeat Purchases in Online Retail}
\newif\ifshowhints
\showhintsfalse      % Hide template instructions in final report
%=============================================================

%-------------------------------------------- fonts (no need to edit)
\usepackage{iftex}
\ifPDFTeX
  \usepackage[T1]{fontenc}
  \usepackage[utf8]{inputenc}
  \IfFileExists{carlito.sty}%
    {\usepackage[sfdefault]{carlito}}%              Calibri look
    {\usepackage[scaled=0.92]{helvet}\renewcommand{\familydefault}{\sfdefault}}%
  \usepackage{courier}
\else
  \usepackage[no-math]{fontspec}
  \setmainfont{Carlito}[Extension=.ttf, UprightFont=*-Regular, BoldFont=*-Bold,
    ItalicFont=*-Italic, BoldItalicFont=*-BoldItalic]
\fi

%----------------------------------------------------------- packages
\usepackage[margin=18mm,bottom=20mm,columnsep=7mm]{geometry}
\usepackage{amsmath}
\usepackage{xcolor}
\usepackage{graphicx}
\usepackage{tabularx}
\usepackage{array}
\usepackage{booktabs}
\usepackage{enumitem}
\usepackage{titlesec}
\usepackage{needspace}
\usepackage[font=small,labelfont=bf]{caption}
\usepackage[hidelinks]{hyperref}
\hypersetup{pdftitle={\projecttitle},
            pdfsubject={CSE558 Data Science, Monsoon 2026: Midterm Project Report}}

%------------------------------------------------------------- layout
\definecolor{eg}{HTML}{7396C4}         % worked example (light blue)
\definecolor{hintgrey}{gray}{0.58}     % instructions (light grey)

\titleformat{\section}{\normalfont\large\bfseries}{\thesection.}{0.5em}{}
\titlespacing*{\section}{0pt}{12pt}{3pt}
\titleformat{\subsection}{\normalfont\normalsize\bfseries}{\thesubsection}{0.5em}{}
\titlespacing*{\subsection}{0pt}{10pt}{2pt}
\newcommand{\sectionbreak}{\needspace{10\baselineskip}}
\newcommand{\subsectionbreak}{\needspace{10\baselineskip}}

\setlength{\parindent}{0pt}
\setlength{\parskip}{4pt}
\setlist{leftmargin=1.2em, itemsep=2pt, topsep=2pt, parsep=0pt}
\renewcommand{\arraystretch}{1.25}
\setlength{\tabcolsep}{4pt}
\newcolumntype{L}[1]{>{\raggedright\arraybackslash}p{#1}}
\newcolumntype{Y}{>{\raggedright\arraybackslash}X}
\newcolumntype{K}[1]{>{\raggedright\arraybackslash\bfseries\boldmath}p{#1}}  % bold labels

%--------------------------------------------------------- shorthands
% Instruction line. Delete these as you fill the template in.
\newcommand{\hint}[1]{\ifshowhints{\footnotesize\itshape\color{hintgrey}#1\par}\fi}
% Worked example: inline or in a table cell. REPLACE the contents.
\newcommand{\eg}[1]{\textcolor{eg}{#1}}
% Worked example running over a whole paragraph.
\newcommand{\egblock}[1]{{\color{eg}#1\par}}
% Bold label followed by an example.
\newcommand{\field}[2]{\textbf{#1}~\eg{#2}\par}
% Part a)-c) of a test, and a small bold heading.
\newcommand{\tpart}[2]{\par\needspace{12\baselineskip}\smallskip\textbf{#1)~#2}\par}
\newcommand{\minihead}[1]{\par\needspace{8\baselineskip}\smallskip\textbf{#1}\par}
% Column and file names.
\newcommand{\code}[1]{\texttt{#1}}
% \figfile{file}{height}: your plot, or a labelled box until the file exists.
\newcommand{\figfile}[2]{%
  \IfFileExists{#1}%
    {\includegraphics[width=\linewidth,height=#2,keepaspectratio]{#1}}%
    {\fbox{\parbox[c][#2][c]{\dimexpr\linewidth-2\fboxsep-2\fboxrule\relax}%
      {\centering\small\itshape\color{hintgrey}Your plot: \detokenize{#1}}}}}

%=======================================================================

\begin{document}
\twocolumn[%
{\small\bfseries CSE558 Data Science\ \ \textbullet\ \ Midterm Project Report\ \ \textbullet\ \ Monsoon 2026\par}
\vspace{4pt}
{\LARGE\bfseries\raggedright\projecttitle\par}
\vspace{6pt}\rule{\linewidth}{1pt}\vspace{4pt}
{\small
\begin{tabularx}{\linewidth}{@{}K{2.4cm} Y@{}}
\toprule
Group name & \textbf{[Enter group name]}\\
Dataset & UCI Online Retail, DOI: 10.24432/C5BW33\\
Code & \textbf{[Insert public repository URL after publishing]}\\
\bottomrule
\end{tabularx}}
\vspace{4pt}
{\small
\begin{tabularx}{\linewidth}{@{}L{0.5cm} L{3.4cm} L{2.3cm} L{4.8cm} Y@{}}
\toprule
\textbf{\#} & \textbf{Name} & \textbf{Roll no.} & \textbf{Institute email} & \textbf{Contribution}\\
\midrule
1 & \textbf{[Your name]} & \textbf{[Roll no.]} & \textbf{[Institute email]} & Dataset, analysis, report and code\\
\bottomrule
\end{tabularx}}
\vspace{8pt}]
\section{Dataset}
\field{Source:}{UCI Machine Learning Repository, Online Retail (Chen, 2015), \url{https://doi.org/10.24432/C5BW33}; CC BY 4.0.}
\field{Raw data:}{541,909 rows and 8 original attributes; 1 December 2010--9 December 2011. Each row is an invoice product-line. The working audit adds one derived cancellation flag, so its intermediate schema has 9 columns; the prepared transaction table has additional derived fields.}
\field{Encoded-size check:}{Country has 38 levels in the encoded-size audit. One-hot encoding replaces the Country column with 38 indicator columns while retaining the other 7 original attributes: $7+38=45$ features. Thus $541{,}909\times45=24{,}385{,}905$ encoded cells, above 3,000,000. If the derived cancellation flag is retained as a feature, there are 46 columns and $541{,}909\times46=24{,}927{,}814$ cells. Confirm the actual category count and encoded matrix shape when running the script; this arithmetic is not a substitute for the runtime shape check.}
{\footnotesize
\begin{tabularx}{\linewidth}{@{}L{2.0cm} L{1.7cm} L{1.4cm} Y@{}}
\toprule
\textbf{Feature} & \textbf{Type} & \textbf{Unit/range} & \textbf{Meaning and missingness}\\
\midrule
\code{InvoiceNo} & Nominal & Invoice ID & Invoice identifier; cancellation invoices often begin with C. No missingness reported in audit.\\
\code{StockCode} & Nominal & Product ID & Product identifier. No missingness reported in audit.\\
\code{Description} & Text & Product text & Product description; 1,454 rows missing.\\
\code{Quantity} & Discrete & Integer units & Units on line; 10,624 rows non-positive.\\
\code{InvoiceDate} & Datetime & 2010-12-01 to 2011-12-09 & Invoice timestamp.\\
\code{UnitPrice} & Continuous & GBP/unit & Unit price; 2,517 rows non-positive.\\
\code{CustomerID} & Nominal ID & Customer ID & Customer identifier; 135,080 rows missing.\\
\code{Country} & Nominal & Country labels & Customer country; encoded as one-hot indicators for the size check.\\
\bottomrule
\end{tabularx}}
\section{Data preparation}
{\footnotesize
\begin{tabularx}{\linewidth}{@{}L{2.1cm} L{2.8cm} Y@{}}
\toprule
\textbf{Issue} & \textbf{Action} & \textbf{Reason / alternative considered}\\
\midrule
Exact duplicate lines (5,268) & Remove exact duplicates & Prevent identical repeated records from inflating quantities and spend; retain non-identical lines.\\
Missing CustomerID (135,080) & Exclude from customer-level analysis; keep in raw audit & Imputing IDs would fabricate customer histories and repeat-purchase labels.\\
Missing descriptions (1,454) & Retain where other transaction fields are usable; description is not needed for customer-level target & Dropping all such rows would unnecessarily discard valid purchases.\\
Cancellation invoices (9,288) & Exclude from positive-sale cohort using invoice cancellation flag & Returns/cancellations should not count as positive purchases; still report them in audit.\\
Quantity $\leq0$ (10,624) & Exclude from positive-purchase analysis & Non-positive quantities represent returns/corrections rather than a positive purchase.\\
UnitPrice $\leq0$ (2,517) & Exclude from positive-purchase analysis & Non-positive prices do not represent positive sales value.\\
Skewed spend / outliers & Retain valid positive values; use log spend as an ML feature & High spend can be real customer behaviour. Automatic capping/deletion may erase genuine high-value customers.\\
Categorical country & One-hot encode in ML pipeline & Avoid imposing a false numeric order on country labels; unknown test categories are ignored by encoder.\\
\bottomrule
\end{tabularx}}
\section{What changed after preparation}
{\footnotesize
\begin{tabularx}{\linewidth}{@{}L{3.3cm} L{2.0cm} Y@{}}
\toprule
\textbf{Measure} & \textbf{Before} & \textbf{After}\\
\midrule
Rows & 541,909 raw lines & 392,692 positive-sale lines\\
Original columns & 8 & Prepared table has derived fields; customer table has 13 fields\\
Customers with IDs & 4,372 & 4,338 customers after filtering\\
Distinct invoices & 25,900 raw audit & 18,532 positive-sale invoices\\
Exact duplicates & 5,268 & 0 after exact duplicate removal\\
Missing CustomerID & 135,080 & Excluded from customer-level table, not imputed\\
Cancellation invoice lines & 9,288 & Excluded from positive-sale cohort\\
Non-positive quantity lines & 10,624 & Excluded from positive-sale cohort\\
Non-positive price lines & 2,517 & Excluded from positive-sale cohort\\
\bottomrule
\end{tabularx}}
The prepared positive-sale table has 392,692 lines, 18,532 distinct invoices and 4,338 identified customers. These row counts are not additive because a row may meet more than one exclusion condition. The final customer table has one row per customer and supports inference without treating multiple invoices from one customer as independent observations.
\section{Basic questions and problem statement}
\begin{figure}[!ht]\figfile{../outputs/eda_purchase_frequency_exploration.png}{3.2cm}
\caption{Purchase frequency in the exploration split is right-skewed: many customers have one invoice, with a long tail of repeat buyers.}\end{figure}
\begin{figure}[!ht]\figfile{../outputs/eda_spend_repeat_exploration.png}{3.2cm}
\caption{Exploration-split repeat customers show a broader total-spend distribution; display is capped at the exploration 99th percentile for readability.}\end{figure}
\begin{figure}[!ht]\figfile{../outputs/eda_monthly_invoices_exploration.png}{3.2cm}
\caption{Monthly invoice volume among exploration customers varies over time; December 2011 is incomplete because records end on 9 December.}\end{figure}
\minihead{Split before exploration}
Customer-level data were split 80/20 with stratification on observed repeat-customer status and seed 558: exploration $n=3,470$ and confirmation $n=868$. EDA figures above use exploration customers only. The three tests below were run on the held-out confirmation customers; the median first-invoice threshold was set from the exploration split only (£300.655).
\minihead{Obvious questions}
\begin{enumerate}
\item How many customers buy once versus repeatedly?
\item Which customers spend the most over the observed period?
\item How does purchase activity vary by month?
\end{enumerate}
\minihead{Problem statements}
\noindent\fbox{\begin{minipage}{\dimexpr\linewidth-2\fboxsep-2\fboxrule\relax}
\textbf{P1 (primary).} Is the probability of repeat purchasing different for customers whose first invoice value is above the exploration-set median compared with those at or below it?

\smallskip
\textbf{P2.} Does customer-level mean invoice value differ between repeat and one-time customers?

\smallskip
\textbf{P3.} Does repeat-purchase probability differ between UK and non-UK customers?
\end{minipage}}
\field{Novelty:}{The obvious questions describe historical volume and spend. These statements test pre-specified customer-level contrasts in a held-out confirmation sample, use a future-relevant retention outcome for P1, and explicitly avoid counting each invoice as an independent customer. They remain observational associations, not causal claims.}
\section{Statistical inferences}
All tests use the confirmation set (868 customers), two-sided alternatives, and $\alpha=0.05$. Practical thresholds below are chosen for interpretation in this report and were not preregistered; therefore practical-significance conclusions are exploratory. For two proportions, the minimum practically meaningful difference is 5 percentage points. For average invoice value, it is £50 per customer.
\subsection*{Test 1 (primary): first-invoice spend and repeat purchasing}
\tpart{a}{Hypotheses and method}
\begin{tabularx}{\linewidth}{@{}K{3.2cm} Y@{}}
\toprule
$H_0$ & $p_{\mathrm{above}}-p_{\mathrm{at/below}}=0$\\
$H_1$ & $p_{\mathrm{above}}-p_{\mathrm{at/below}}\ne0$\\
Significance & $\alpha=0.05$; two-sided\\
Test & Two-proportion z-test\\
Practical threshold & 5 percentage points\\
\bottomrule
\end{tabularx}
The test compares independent customers in the two first-invoice groups. A two-proportion z-test is suitable for the difference in proportions at these sample sizes.
\tpart{b}{Other tests considered}
A chi-square test of independence on the corresponding $2\times2$ table is asymptotically equivalent and would give the same two-sided conclusion. Fisher's exact test avoids the large-sample approximation but is not necessary with these counts. A one-sided test was not used because the question allows either direction.
\tpart{c}{Result}
Above-median group: 290/424 = 68.40\%; at/below-median: 279/444 = 62.84\%. Difference = 5.56 percentage points; 95\% CI [-0.75, 11.87] percentage points; $z=1.723$, degrees of freedom: not applicable for a z-test, $p=0.085$. Fail to reject $H_0$. The point estimate is just above the 5-point practical threshold, but the interval includes zero and values below the threshold; evidence is inconclusive.
\subsection*{Test 2: customer-level average invoice value by repeat status}
\tpart{a}{Hypotheses and method}
\begin{tabularx}{\linewidth}{@{}K{3.2cm} Y@{}}
\toprule
$H_0$ & $\mu_{\mathrm{repeat}}-\mu_{\mathrm{one-time}}=0$\\
$H_1$ & $\mu_{\mathrm{repeat}}-\mu_{\mathrm{one-time}}\ne0$\\
Significance & $\alpha=0.05$; two-sided\\
Test & Welch two-sample t-test\\
Practical threshold & £50 per customer\\
\bottomrule
\end{tabularx}
Each customer contributes one average invoice value. Welch's test does not require equal population variances.
\tpart{b}{Other tests considered}
The pooled t-test assumes equal variances; we avoid that assumption. A Mann--Whitney test compares ranks/distributions, not specifically means. A bootstrap confidence interval would be a useful robustness check for the skewed values. A test on invoice-level rows was rejected because repeat customers contribute multiple invoices, creating pseudoreplication.
\tpart{c}{Result}
Repeat customers: mean £398.95 ($n=569$); one-time customers: £384.08 ($n=299$). Difference = £14.88; 95\% CI [-£54.48, £84.23]; $t=0.421$, $df=530.68$, $p=0.674$. Fail to reject $H_0$. The point estimate is below £50, but the interval includes differences larger than the practical threshold in either direction; the data do not establish equivalence.
\subsection*{Test 3: UK versus non-UK repeat-purchase rates}
\tpart{a}{Hypotheses and method}
\begin{tabularx}{\linewidth}{@{}K{3.2cm} Y@{}}
\toprule
$H_0$ & $p_{\mathrm{UK}}-p_{\mathrm{nonUK}}=0$\\
$H_1$ & $p_{\mathrm{UK}}-p_{\mathrm{nonUK}}\ne0$\\
Significance & $\alpha=0.05$; two-sided\\
Test & Two-proportion z-test\\
Practical threshold & 5 percentage points\\
\bottomrule
\end{tabularx}
\tpart{b}{Other tests considered}
A chi-square test of independence is asymptotically equivalent. Fisher's exact test is a possible small-cell alternative, although the non-UK group has 95 customers and observed success/failure counts are adequate for the normal approximation. A logistic regression adjusted for country/customer covariates could address confounding but answers a different, adjusted question.
\tpart{c}{Result}
UK repeat rate: 508/773 = 65.72\%; non-UK: 61/95 = 64.21\%. Difference = 1.51 percentage points; 95\% CI [-8.70, 11.71] percentage points; $z=0.292$, degrees of freedom: not applicable for a z-test, $p=0.770$. Fail to reject $H_0$. The interval is wide and does not rule out practically meaningful differences.
\section{Scalable machine learning}
The ML target is whether a customer with purchase history before 1 September 2011 purchases during the next 90 days. Features are calculated only from the history period. The cohort contains 3,317 customers; 1,864 (56.2\%) have a positive label. On a fixed 830-customer test set, training size was varied across 621, 1,243 and 2,487 customers using the same seed, preprocessing and model settings.
{\footnotesize
\begin{tabularx}{\linewidth}{@{}L{2.7cm} r r r r@{}}
\toprule
\textbf{Model} & \textbf{Train n} & \textbf{F1} & \textbf{ROC-AUC} & \textbf{Fit (s)}\\
\midrule
Logistic regression & 621 & 0.671 & 0.718 & 0.016\\
Logistic regression & 1,243 & 0.686 & 0.726 & 0.016\\
Logistic regression & 2,487 & 0.693 & 0.732 & 0.027\\
Shallow tree & 621 & 0.636 & 0.701 & 0.009\\
Shallow tree & 1,243 & 0.667 & 0.717 & 0.032\\
Shallow tree & 2,487 & 0.623 & 0.727 & 0.039\\
\bottomrule
\end{tabularx}}
Logistic regression achieves the best full-sample F1 and ROC-AUC. Tree performance is not monotonic with training size. Fit times are machine-specific and this small customer-level dataset does not prove large-scale scalability; the experiment is a controlled sample-size comparison. The evaluation split is stratified randomly within a temporal cohort, not a pure later-period holdout, which limits deployment claims.
\section{Limitations}
\begin{itemize}
\item \textbf{Coverage.} One UK retailer over roughly one year; December 2011 is incomplete.
\item \textbf{Missing customer IDs.} 135,080 raw rows cannot be reliably linked to customer histories and are excluded from customer-level analyses.
\item \textbf{Practical thresholds.} Thresholds were selected for interpretation in this report, not preregistered before analysis.
\item \textbf{Association, not causation.} Spending, geography and repeat status may be related to unobserved customer characteristics.
\item \textbf{Encoding verification.} The size arithmetic uses 38 observed Country levels from the prior audit; the exact encoded matrix shape must be verified by running the included script on the original raw file.
\item \textbf{Submission metadata.} Add the actual group name, name, roll number, institute email and public code URL before submission.
\end{itemize}
\end{document}
